\documentclass{article}
\usepackage[margin=1.1in]{geometry}
\usepackage{graphicx}
\usepackage{amsmath}
\usepackage{amssymb}
\usepackage{hyperref}
\usepackage{booktabs}
\usepackage{xcolor}
\usepackage{listings}
\usepackage{tcolorbox}
\usepackage{pdfpages}
\tcbuselibrary{listings, skins, breakable}

% Output box environment (black background, white text)
\newtcolorbox{outputbox}[1]{
    arc=3mm,
    boxrule=1pt,
    colback=black!85!white,
    colframe=black,
    coltext=white,
    title={\textbf{#1}},
    fontupper=\ttfamily\scriptsize,
    left=2mm, right=2mm, top=2mm, bottom=2mm,
    fontupper=\ttfamily\scriptsize\obeylines,
    breakable
}

% Python code box environment
\lstdefinestyle{pythonstyle}{
    language=Python,
    basicstyle=\ttfamily\scriptsize,
    keywordstyle=\color{blue!80!black}\bfseries,
    stringstyle=\color{green!60!black},
    commentstyle=\color{gray}\itshape,
    showstringspaces=false,
    numbers=left,
    numberstyle=\tiny\color{gray},
    numbersep=3.5pt,
    breaklines=true,
    tabsize=4
}
\newtcblisting{pythonbox}[1]{
    arc=3mm,
    boxrule=1pt,
    colback=yellow!5!white,
    colframe=yellow!50!black,
    title={\textbf{#1}},
    listing only,
    listing style=pythonstyle,
    breakable
}

% ── Title ─────────────────────────────────────────────────────────────────────
\begin{document}

\begin{titlepage}
\begin{center}


\includegraphics[width=0.55\textwidth]{images/ozu.png}\\[2cm]

\rule{\linewidth}{0.5mm}\\[0.4cm]

{ \LARGE
  \textbf{PHYS551: Advanced Monte Carlo Methods}\\[0.4cm]
  \emph{Homework 4}\\[0.4cm]
}
\rule{\linewidth}{0.5mm}\\[0.4cm]

\hfill\hfill

\linewidth=0.8\textwidth
\begin{minipage}{0.4\textwidth}
    \begin{flushleft} \large
        \emph{Student:}\\
        Ömer Faruk \textsc{Avcı}
    \end{flushleft}
\end{minipage}
~
\begin{minipage}{0.5\textwidth}
    \begin{flushright} \large
        \emph{Instructor:}\\
        Prof. Dr. Taylan \textsc{Akdo\u{g}an}\\[0.1cm]
    \end{flushright}
\end{minipage}

\vfill

\textsc{\large Department of Electrical and Electronics Engineering}\\[0.4cm]

{\large \today}

\end{center}
\end{titlepage}

% ── Disclosures ───────────────────────────────────────────────────────────────
\newpage
\section*{Disclosures and References}

\textbf{AI Usage Disclosure:} In accordance with the course policy, I would like to declare the extent of AI assistance used in this assignment:

\begin{itemize}
    \item \textbf{Analytical Derivations:} All mathematical derivations were performed manually without the use of AI tools.
    \item \textbf{Code Implementation:} The code was developed by me. AI tools were consulted to verify correctness and refine code comments.
    \item \textbf{Report Writing:} The explanations and discussions were written by me. AI tools were used to check for grammatical errors and to give feedback on the overall structure.
\end{itemize}

\newpage

% =============================================================
% PROBLEM 1
% =============================================================
\includepdf{problem_1_derivation}
\section*{Problem 2: Mixing Time and The Spectral Gap (Computational)}

\subsection*{Part (a): Simulating the Evolution}

Starting from $\mu^{(0)} = [1,\, 0,\, 0]$, we compute $\mu^{(t)} = \mu^{(0)} P^t$ by iteratively multiplying by $P$.

\begin{pythonbox}{Python Source Code for Problem 2}
"""

# PHYS 551: Advanced Monte Carlo Methods
# Homework 3 - Source Code
# Student: Omer Faruk Avci
# Instructor: Prof. Dr. Taylan Akdogan

"""
import numpy as np
import matplotlib.pyplot as plt

"""Problem 2: Mixing Time and The Spectral Gap (Computational)"""


# Transition matrix
P = np.array([
    [1/2, 1/2,   0  ],
    [1/5, 1/2, 3/10 ],
    [0,   2/5,  3/5 ]
])

# Initial distribution given in the problem
mu_init = np.array([1.0, 0.0, 0.0])

t_max = 40
mu_values = np.zeros((t_max + 1, 3))
mu_values[0] = mu_init

for t in range(1, t_max + 1):
    mu_values[t] = mu_values[t-1] @ P  # mu^(t) = mu^(0) * P^t

# Part b: Plot
pi = np.array([8/43, 20/43, 15/43])
time = np.arange(t_max + 1)

plt.figure(figsize=(10, 6))
colors = ['red', 'orange', 'blue']
for i, c in enumerate(colors):
    plt.plot(time, mu_values[:, i], f'-o', color=c, label=f'State {i+1}')
    plt.axhline(pi[i], color=c, linestyle='--', alpha=0.6)

plt.xlabel('Time')
plt.ylabel('Probability')
plt.title('Markov Chain Evolution')
plt.legend(); plt.grid(True, alpha=0.6)
plt.tight_layout()
plt.savefig("markov_evolution.png")

# Part c: Eigenvalues
eigenvalues = np.linalg.eigvals(P)
eigenvalues_sorted = sorted(eigenvalues, key=abs, reverse=True)

lambda_1 = eigenvalues_sorted[0]
lambda_2 = eigenvalues_sorted[1]
gamma    = 1 - abs(lambda_2)
tau      = -1 / np.log(abs(lambda_2))

print(f"lambda1 = {lambda_1:.6f}")
print(f"lambda2 = {lambda_2:.6f}")
print(f"Spectral Gap = {gamma:.6f}")
print(f"Relaxation Time = {tau:.6f}")
\end{pythonbox}

\subsection*{Part (b): Visualization}

The figure below shows $\mu^{(t)}$ for each state as a function of $t \in [0, 40]$. The horizontal dashed lines represent the exact analytical stationary values found in Problem 1b.

\begin{figure}[h!]
    \centering
    \includegraphics[width=0.88\textwidth]{images/problem_2_figure.png}
    \caption{Markov Chain Evolution}
\end{figure}

The chain converges visually to the stationary distribution by approximately $t \approx 8$--$10$, consistent with the spectral gap analysis below.

\newpage
\subsection*{Part (c): Spectral Gap and Relaxation Time}

The three eigenvalues of $P$, sorted by absolute value, are:

\begin{table}[h!]
\centering
\renewcommand{\arraystretch}{1.3}
\begin{tabular}{|l|c|}
\hline
\textbf{Quantity} & \textbf{Value} \\ \hline
Largest eigenvalue $\lambda_1$ & $1.000000$ \\ \hline
Second largest eigenvalue $|\lambda_2|$ & $0.544949$ \\ \hline
Spectral Gap $\gamma = 1 - |\lambda_2|$ & $0.455051$ \\ \hline
Relaxation time $\tau = -1/\ln|\lambda_2|$ & $1.647275$ \\ \hline
\end{tabular}
\caption{Eigenvalue analysis and derived mixing quantities.}
\end{table}

\begin{outputbox}{Terminal Output}
lambda 1 = 1.000000
lambda 2 = 0.544949
Spectral Gap = 0.455051
Relaxation Time = 1.647275
\end{outputbox}

\textbf{Interpretation of $\tau$:}
The total variation distance between $\mu^{(t)}$ and $\pi$ decays approximately as:
\[
\|\mu^{(t)} - \pi\|_{\mathrm{TV}} \propto e^{-t/\tau}
\]
At $t = \tau \approx 1.65$ steps, the distance has already decreased by a factor of $1/e$. In practice, convergence is accepted when $t \approx 5\tau$:
\[
5\tau \approx 8.24 \text{ steps}
\]
This is fully consistent with the plot in Part (b), where all three state probabilities have settled onto their stationary dashed lines by $t \approx 8$--$10$.

\end{document}